\documentclass[11pt]{article}
\usepackage{amsmath,amssymb,amsthm}
\usepackage{mathpartir}
\usepackage[margin=2.5cm]{geometry}
\usepackage{hyperref}

\newcommand{\U}{\mathsf{U}}
\newcommand{\Prop}{\mathsf{Prop}}
\newcommand{\Pit}[2]{\Pi(#1)\,#2}
\newcommand{\lam}[2]{\lambda(#1)\,#2}
\newcommand{\ctx}[1]{#1\ \mathsf{ctx}}
\newcommand{\ty}[3]{#1 \vdash #2 : #3}
\newcommand{\cv}[4]{#1 \vdash #2 = #3 : #4}
\newcommand{\hr}{\mathrel{\to^{*}_{\mathsf{h}}}}
\newcommand{\hrone}{\mathrel{\to_{\mathsf{h}}}}
\newcommand{\code}[1]{\mathsf{#1}}
\newcommand{\IsProp}{\mathsf{IsProp}}
\newcommand{\agda}[1]{\texttt{#1}}

\theoremstyle{plain}
\newtheorem{theorem}{Theorem}
\newtheorem{lemma}[theorem]{Lemma}
\newtheorem{corollary}[theorem]{Corollary}
\theoremstyle{remark}
\newtheorem*{remark}{Remark}

\title{Dependent type theory with $\U : \U$ and an impredicative\\
  universe $\Prop \subseteq \U$ of definitionally proof-irrelevant propositions:\\
  typing rules and main results}
\author{Formalisation in \agda{MINPROP/} (Agda 2.9.0, \agda{--without-K})}
\date{21 October 2026}

\begin{document}
\maketitle

\begin{abstract}
We present a minimal dependent type theory with $\Pi$-types, a universe
$\U$ with $\U : \U$, and a universe $\Prop : \U$ of propositions with
$\Prop \subseteq \U$, closed under impredicative $\Pi$, and with
definitional proof irrelevance. Using a domain-theoretic model and a
logical relation, we prove that $\Pi$ is injective and that typing is
preserved by head reduction. Everything is formalised in Agda with no
postulates, no holes and no termination or positivity pragmas.
\end{abstract}

\section{Syntax}

Terms are given in de~Bruijn style (\agda{Syntax/Raw.agda}):
\[
  M, N, A, B ::= x \mid \U \mid \Prop \mid \Pit{x:A}{B} \mid \lam{x:A}{M} \mid M\,N .
\]
Contexts are snoc-lists $\Gamma ::= \cdot \mid \Gamma, x{:}A$. The file
\agda{Syntax/Typing.agda} defines three mutually inductive judgements:
\[
  \ctx{\Gamma} \qquad \ty{\Gamma}{M}{A} \qquad \cv{\Gamma}{M}{N}{A}
\]

Head reduction (\agda{Syntax/Reduction.agda}) is generated by
\[
  (\lam{x:A}{M})\,N \hrone M[N/x]
  \qquad\qquad
  \inferrule{M \hrone M'}{M\,N \hrone M'\,N}
\]
and $\hr$ denotes its reflexive--transitive closure. In particular $\U$,
$\Prop$ and $\Pi$-types are head normal.

\section{Typing rules}

\subsection*{Contexts}
\begin{mathpar}
  \inferrule{ }{\ctx{\cdot}}
  \and
  \inferrule{\ty{\Gamma}{A}{\U}}{\ctx{\Gamma, x{:}A}}
\end{mathpar}

\subsection*{Typing}
\begin{mathpar}
  \inferrule*[right=Var]{\ctx{\Gamma} \\ (x{:}A) \in \Gamma}{\ty{\Gamma}{x}{A}}
  \and
  \inferrule*[right=Conv]{\ty{\Gamma}{M}{A} \\ \cv{\Gamma}{A}{B}{\U} \\ \ty{\Gamma}{B}{\U}}
                         {\ty{\Gamma}{M}{B}}
  \\
  \inferrule*[right=U]{\ctx{\Gamma}}{\ty{\Gamma}{\U}{\U}}
  \and
  \inferrule*[right=Pi]{\ty{\Gamma}{A}{\U} \\ \ty{\Gamma, x{:}A}{B}{\U}}
                       {\ty{\Gamma}{\Pit{x:A}{B}}{\U}}
  \\
  \inferrule*[right=Lam]{\ty{\Gamma}{A}{\U} \\ \ty{\Gamma, x{:}A}{B}{\U} \\ \ty{\Gamma, x{:}A}{M}{B}}
                        {\ty{\Gamma}{\lam{x:A}{M}}{\Pit{x:A}{B}}}
  \and
  \inferrule*[right=App]{\ty{\Gamma}{A}{\U} \\ \ty{\Gamma, x{:}A}{B}{\U} \\
                         \ty{\Gamma}{M}{\Pit{x:A}{B}} \\ \ty{\Gamma}{N}{A}}
                        {\ty{\Gamma}{M\,N}{B[N/x]}}
\end{mathpar}

\paragraph{The universe of propositions.}
\begin{mathpar}
  \inferrule*[right=Prop]{\ctx{\Gamma}}{\ty{\Gamma}{\Prop}{\U}}
  \and
  \inferrule*[right=Prop-U]{\ty{\Gamma}{A}{\Prop}}{\ty{\Gamma}{A}{\U}}
  \and
  \inferrule*[right=Pi-Prop]{\ty{\Gamma}{A}{\U} \\ \ty{\Gamma, x{:}A}{B}{\Prop}}
                            {\ty{\Gamma}{\Pit{x:A}{B}}{\Prop}}
\end{mathpar}
The domain $A$ in \textsc{Pi-Prop} is any type, for instance $\Prop$ itself,
so $\Prop$ is impredicative: $\ty{\Gamma}{\Pit{X:\Prop}{X}}{\Prop}$.
There is no rule $\Prop : \Prop$.

\subsection*{Conversion}
\begin{mathpar}
  \inferrule*[right=Refl]{\ty{\Gamma}{M}{A}}{\cv{\Gamma}{M}{M}{A}}
  \and
  \inferrule*[right=Sym]{\cv{\Gamma}{M}{N}{A}}{\cv{\Gamma}{N}{M}{A}}
  \and
  \inferrule*[right=Trans]{\cv{\Gamma}{M}{N}{A} \\ \cv{\Gamma}{N}{P}{A}}{\cv{\Gamma}{M}{P}{A}}
  \\
  \inferrule*[right=Conv-Eq]{\cv{\Gamma}{M}{N}{A} \\ \cv{\Gamma}{A}{B}{\U} \\ \ty{\Gamma}{B}{\U}}
                            {\cv{\Gamma}{M}{N}{B}}
  \and
  \inferrule*[right=Beta]{\ty{\Gamma}{A}{\U} \\ \ty{\Gamma, x{:}A}{B}{\U} \\
                          \ty{\Gamma, x{:}A}{M}{B} \\ \ty{\Gamma}{N}{A}}
                         {\cv{\Gamma}{(\lam{x:A}{M})\,N}{M[N/x]}{B[N/x]}}
  \\
  \inferrule*[right=Pi-Cong]{\ty{\Gamma}{A}{\U} \\ \ty{\Gamma, x{:}A}{B}{\U} \\ \ty{\Gamma, x{:}A}{B'}{\U} \\\\
                             \cv{\Gamma}{A}{A'}{\U} \\ \cv{\Gamma, x{:}A}{B}{B'}{\U}}
                            {\cv{\Gamma}{\Pit{x:A}{B}}{\Pit{x:A'}{B'}}{\U}}
  \\
  \inferrule*[right=Funext]{\ty{\Gamma}{A}{\U} \\ \cv{\Gamma, x{:}A}{M\,x}{N\,x}{B} \\
                            \ty{\Gamma}{M}{\Pit{x:A}{B}} \\ \ty{\Gamma}{N}{\Pit{x:A}{B}}}
                           {\cv{\Gamma}{M}{N}{\Pit{x:A}{B}}}
  \\
  \inferrule*[right=App-Fun]{\ty{\Gamma}{A}{\U} \\ \ty{\Gamma, x{:}A}{B}{\U} \\
                             \cv{\Gamma}{M}{M'}{\Pit{x:A}{B}} \\ \ty{\Gamma}{N}{A}}
                            {\cv{\Gamma}{M\,N}{M'\,N}{B[N/x]}}
  \and
  \inferrule*[right=App-Arg]{\ty{\Gamma}{A}{\U} \\ \ty{\Gamma, x{:}A}{B}{\U} \\
                             \ty{\Gamma}{M}{\Pit{x:A}{B}} \\ \cv{\Gamma}{N}{N'}{A}}
                            {\cv{\Gamma}{M\,N}{M\,N'}{B[N/x]}}
\end{mathpar}
In \textsc{Funext}, $M$ and $N$ are weakened in the premise.

\paragraph{Conversion rules for $\Prop$.}
\begin{mathpar}
  \inferrule*[right=Proof-Irr]{\ty{\Gamma}{A}{\Prop} \\ \ty{\Gamma}{M}{A} \\ \ty{\Gamma}{N}{A}}
                              {\cv{\Gamma}{M}{N}{A}}
  \and
  \inferrule*[right=Prop-U-Eq]{\cv{\Gamma}{M}{N}{\Prop}}{\cv{\Gamma}{M}{N}{\U}}
  \and
  \inferrule*[right=U-Prop-Eq]{\ty{\Gamma}{A}{\Prop} \\ \ty{\Gamma}{B}{\Prop} \\ \cv{\Gamma}{A}{B}{\U}}
                               {\cv{\Gamma}{A}{B}{\Prop}}
  \\
  \inferrule*[right=Pi-Prop-Cong]{\ty{\Gamma}{A}{\U} \\ \ty{\Gamma, x{:}A}{B}{\Prop} \\ \ty{\Gamma, x{:}A}{B'}{\Prop} \\\\
                                  \cv{\Gamma}{A}{A'}{\U} \\ \cv{\Gamma, x{:}A}{B}{B'}{\Prop}}
                                 {\cv{\Gamma}{\Pit{x:A}{B}}{\Pit{x:A'}{B'}}{\Prop}}
\end{mathpar}

Together, \textsc{Prop-U-Eq} and \textsc{U-Prop-Eq} say that two
propositions are convertible as elements of $\Prop$ exactly when they are
convertible as elements of $\U$.

\begin{remark}
The system contains $\U : \U$ and so is logically inconsistent (Girard's
paradox). The results below do not depend on consistency. They concern the
equational theory: they show that \textsc{Proof-Irr} does not make the
theory degenerate.
\end{remark}

\section{Main results}

\begin{theorem}[$\Pi$-injectivity; \agda{PiInjectivity.piInjectivity}]
If $\cv{\Gamma}{\Pit{x:A}{B}}{\Pit{x:A'}{B'}}{\U}$, then
$\cv{\Gamma}{A}{A'}{\U}$ and $\cv{\Gamma, x{:}A}{B}{B'}{\U}$.
\end{theorem}

\noindent This is derived from the more general statement:

\begin{theorem}[\agda{PiInjectivity.piConv}]
If $\cv{\Gamma}{T}{\Pit{x:A'}{B'}}{\U}$, then there are $A, B$ with
$T \hr \Pit{x:A}{B}$, $\cv{\Gamma}{A}{A'}{\U}$ and
$\cv{\Gamma, x{:}A}{B}{B'}{\U}$.
\end{theorem}

\begin{corollary}
If $\cv{\Gamma}{\Pit{x:A}{B}}{\Pit{x:A'}{B'}}{\Prop}$, then
$\cv{\Gamma}{A}{A'}{\U}$ and $\cv{\Gamma, x{:}A}{B}{B'}{\U}$
(by \textsc{Prop-U-Eq} and Theorem~1).
\end{corollary}

\begin{theorem}[Subject reduction; \agda{SubjectReduction.subject-red1}]
If $\ty{\Gamma}{M}{A}$ and $M \hrone N$, then $\ty{\Gamma}{N}{A}$.
\end{theorem}
In the $\beta$ case, a $\lambda$ typed through \textsc{Prop-U} would need
$\cv{\Gamma}{\U}{\Pit{x:A}{B}}{\U}$. This is refuted by Theorem~2, because
$\U$ is head normal.

\begin{theorem}[No confusion; \agda{NoConfusion}]
In every context, none of the following holds:
$\cv{\Gamma}{\U}{\Prop}{\U}$, $\cv{\Gamma}{\U}{\Pit{x:A}{B}}{\U}$,
$\cv{\Gamma}{\Prop}{\Pit{x:A}{B}}{\U}$.
\end{theorem}

\section{Outline of the proof}

\paragraph{Domain model.}
Types and terms are interpreted by relations $M \Downarrow_\rho u$ to
finite elements $u$ of a domain (\agda{Model/Eval.agda}). The finite
elements are
\[
  u ::= \bot \mid \code{U} \mid \code{Prop} \mid \code{Fun}(g) \mid \Pi(a, f),
\]
where $g, f$ are finite lists of steps. The elements carry an order $\le$,
a compatibility relation and a sup. Membership $u : a$ (``$u$ is an element of
the type $a$'') is defined by structural recursion on a stage index, which
is why no termination pragma is needed. The codes of propositions are given
by a structural, stage-free predicate:
\[
  \IsProp(\bot), \qquad
  \IsProp(\Pi(a, f)) \iff \text{every value of } f \text{ satisfies } \IsProp,
\]
and $u : \code{Prop}$ is defined as $u : \code{U} \wedge \IsProp(u)$. Thus
$\Prop \subseteq \U$ holds semantically by projection.

\begin{lemma}[Collapse; \agda{Domain/Collapse.prop-collapse}]
If $u : a$ and $\IsProp(a)$, then $u = \bot$.
\end{lemma}
\begin{proof}[Proof sketch]
By induction on the rank of $a$. If $a = \bot$, its only element is $\bot$.
If $a = \Pi(b, f)$, an element is either $\bot$ or a coherent $\code{Fun}(g)$
with $g$ non-empty. The first step of $g$ has a value lying in a Prop code of
smaller rank, so that value is $\bot$ by induction. But coherence requires
step values to be different from $\bot$, a contradiction.
\end{proof}

\paragraph{Soundness} (\agda{Model/Soundness.agda}, \agda{Model/SoundnessProp.agda}).
If $\ty{\Gamma}{M}{A}$, then every approximant of $M$ lies below an element
of an approximant of $A$. If $\cv{\Gamma}{M}{N}{A}$, then $M$ and $N$ have the
same approximants. The case \textsc{Proof-Irr} follows from the Collapse
lemma: every approximant of a proof is $\bot$, and $\bot$ approximates
every term.

\paragraph{Logical relation and adequacy}
(\agda{Validity/}, \agda{Adequacy/}). The relation is defined by
stratification on rank:
\begin{itemize}
\item $\mathsf{Val}(M, A, u, a)$ and $\mathsf{EqVal}(M, N, A, u, a)$ for terms;
\item $\mathsf{ValTy}(M, a)$ and $\mathsf{EqValTy}(M, N, a)$ for types.
\end{itemize}
Each clause bundles head reductions together with conversions. At
$\code{Prop}$, $\mathsf{ValTy}(M, \code{Prop})$ says $M \hr \Prop$ and
$\cv{\Gamma}{M}{\Prop}{\U}$, by analogy with $\code{U}$.

The fundamental theorem (\agda{Adequacy/Value.agda}) proves validity of every
derivation, by induction on the derivation. The Prop rules are handled in
\agda{Adequacy/Prop.agda}:
\begin{itemize}
\item \textsc{Pi-Prop} reuses the $\Pi$-into-$\U$ case after weakening the
  codomain premise by \textsc{Prop-U}.
\item \textsc{Prop-U} uses soundness: an approximant of $A : \Prop$ is a
  Prop code.
\item \textsc{U-Prop-Eq}: for a Prop code $u$, the relation at
  $(u, \code{U})$ gives the relation at $(u, \code{Prop})$, since they
  differ only in the component $\mathsf{ValTy}(\U,\code{U})$ versus
  $\mathsf{ValTy}(\Prop,\code{Prop})$.
\item \textsc{Proof-Irr} uses the Collapse lemma: the relation at $u = \bot$
  holds trivially.
\end{itemize}
Theorem~2 is obtained by applying adequacy to the given conversion, at the
environment of all $\bot$s and the identity substitution, at the code
$\Pi(\bot, \mathsf{nil})$.

\paragraph{Formalisation.}
The development consists of 58 Agda files. They are checked with
Agda~2.9.0 and the option \agda{--without-K}; most files also use
\agda{--exact-split}.
It contains no postulates, no holes, and no \agda{TERMINATING},
\agda{NO\_POSITIVITY\_CHECK} or \agda{REWRITE} pragmas. Checking the whole
development takes about 15 seconds. The main files are:
\begin{center}\small
\begin{tabular}{p{9cm}l}
\agda{Syntax/Typing.agda} & the rules above \\
\agda{Domain/MemStage.agda}, \agda{Domain/PropCodes.agda} & membership, $\IsProp$ \\
\agda{Domain/Collapse.agda} & Collapse lemma \\
\agda{Model/Soundness.agda}, \agda{Model/SoundnessProp.agda} & soundness \\
\agda{Validity/Stratified.agda} & the logical relation \\
\agda{Adequacy/Value.agda}, \agda{Adequacy/Prop.agda} & fundamental theorem \\
\agda{PiInjectivity.agda}, \agda{SubjectReduction.agda}, \agda{NoConfusion.agda} & main results
\end{tabular}
\end{center}

\end{document}
